\documentclass[conference]{IEEEtran}
\usepackage{cite}
\usepackage{amsmath,amssymb}
\usepackage{url}

\begin{document}
\title{A Fixed-Mask Pixelated Illumination Pilot Toward Hotspot-Aware Lithographic Optimization}

\author{\IEEEauthorblockN{1\textsuperscript{st} William de Almeida Pavinato}
\IEEEauthorblockA{\textit{Centro de Desenvolvimento Tecnológico}\\
\textit{Universidade Federal de Pelotas}\\
Pelotas, Brasil\\
wapavinato@inf.ufpel.edu.br}
\and
\IEEEauthorblockN{2\textsuperscript{nd} Daniel Henrique S. P. Barretos}
\IEEEauthorblockA{\textit{Centro de Desenvolvimento Tecnológico}\\
\textit{Universidade Federal de Pelotas}\\
Pelotas, Brasil\\
dhspbarretos@inf.ufpel.edu.br}
\and
\IEEEauthorblockN{3\textsuperscript{rd} Luciano Volcan Agostini}
\IEEEauthorblockA{\textit{Centro de Desenvolvimento Tecnológico}\\
\textit{Universidade Federal de Pelotas}\\
Pelotas, Brasil\\
\texttt{agostini@inf.ufpel.edu.b}}
}

\maketitle

\begin{abstract}
This paper documents an implemented, source-only lithography pilot and separates it from a proposed future joint source--mask optimization (SMO) and Neural-ILT system. The pilot optimizes 49 nonnegative pixelated illumination weights with unit total flux while holding seven FIT layouts fixed. It uses a scalar Abbe fixed-basis model, a dose-weighted sigmoid resist, and a static equal-layout soft objective. The hard process-variation-band (PV-band) area is measured separately from that differentiable objective. In the frozen October 9 run, five registered seeds used distinct feasible starts without fallback and completed their 204-step schedules, but all reached the same final FIT source and none qualified under the registered FIT gates. Calibration remained closed, and no final3 layout was indexed or evaluated. A subsequent schema-7 static new-three objective ablation is implemented but has not been benchmarked. These results do not demonstrate a process-window improvement, hotspot detection, generalization, or an official-benchmark evaluation. We report the experiment's scope and measured FIT outcome, identify its limitations, and outline the prospective validation needed before evaluating a joint SMO--Neural-ILT approach.
\end{abstract}

\begin{IEEEkeywords}
computational lithography, illumination-source optimization, source-mask optimization, inverse lithography, process variation
\end{IEEEkeywords}

\section{Introduction}
Optical proximity effects make it difficult to transfer a target layout to a wafer with low print error. Computational lithography studies models and optimization methods for this task, including source optimization and inverse lithography \cite{hopkins,neuralilt2}. Public benchmarks such as LithoBench support comparison across lithography and mask-optimization tasks \cite{lithobench}. Such prior work motivates joint approaches, but does not establish that a particular source-only experiment improves printability.

This paper reports a deliberately narrower result: a fixed-mask pilot that varies only the 49 weights of a pixelated illumination source. All seven FIT layouts remain fixed. No mask pixels or neural-network parameters are optimized, and the experiment does not train a hotspot classifier. An earlier October 7 attempt fell back to the same reference for all five seeds after its jittered starts failed the guard audit. The frozen October 9 schema-6 attempt used five distinct feasible starts, but none produced a checkpoint that met the qualification gate.

This pilot is motivated by earlier work in the broader project on Neural-ILT with PV-band heatmap fine-tuning. Joint SMO--Neural-ILT and heatmap-based hotspot analysis remain future work.

\medskip
\noindent\textbf{Contributions.} This draft:
\begin{itemize}
    \item specifies the implemented 49-weight, fixed-mask illumination experiment and distinguishes its differentiable training objective from the hard PV-band metric;
    \item records the frozen schema-6 outcome, including distinct feasible starts, zero qualified checkpoints, and the closed calibration boundary; and
    \item reports an observed training-surrogate tradeoff that motivates a schema-7 static objective ablation, while separating this unrun ablation and future joint mask/source optimization from measured results.
\end{itemize}

\section{Background and Metrics}
\subsection{Forward Lithography and the Implemented Model}
Wave-optics models describe the formation of lithographic images under partially coherent illumination; Hopkins' diffraction theory is a foundational treatment \cite{hopkins}. The pilot described here uses a scalar Abbe implementation with a precomputed, fixed source basis. For a fixed mask and basis images $B_j$, its aerial intensity is linear in the source weights:
\begin{equation}
 I_{\mathbf{s}}(x,y)=\sum_{j=1}^{49}s_j B_j(x,y),\qquad
 s_j\geq 0,\qquad \sum_{j=1}^{49}s_j=1.
 \label{eq:source}
\end{equation}
This equation describes the implemented source parameterization; it is not a claim that the complete optical and resist pipeline is linear.

At dose $d$, the implementation applies a sigmoid resist response with steepness $\alpha=50$ and threshold $I_{\mathrm{th}}=0.225$, then thresholds the response at 0.5 to obtain a binary print $Z_{d,\mathbf{s}}$. The dose set is $\mathcal{D}=\{0.98,1.00,1.02\}$, and focus is fixed at zero.

\subsection{Hard PV-Band and Soft Objective}
For a layout target $T$, the reported hard PV-band area is the number of pixels printed at some but not all tested doses:
\begin{equation}
 A_{\mathrm{PV}} =
 \left|\left(\bigcup_{d\in\mathcal{D}}Z_{d,\mathbf{s}}\right)
 \setminus
 \left(\bigcap_{d\in\mathcal{D}}Z_{d,\mathbf{s}}\right)\right|.
 \label{eq:pv}
\end{equation}
Nominal $L_2$ is the number of binary mismatches to $T$ at dose 1.00. Worst-dose $L_2$ is the maximum mismatch count over the three doses. These are pixel counts on the evaluated raster.

The reported schema-6 optimization uses the existing differentiable, target-aware critical-corner softcount, averaged with static equal weight over the seven FIT layouts. This soft objective and its associated heatmap are optimization diagnostics; they are not the hard PV-band area in (\ref{eq:pv}), a labeled hotspot map, or a hotspot-classification score. In particular, no hotspot precision, recall, or F1 score is measured in this pilot. The separate schema-7 ablation changes the static soft-loss average as described in Section~\ref{sec:limits}; it has not been run.

\section{Implemented Pilot}
\subsection{Source and Layout Scope}
The pupil is discretized on a $9\times9$ grid in pupil coordinates $(u,v)$. The 49 permitted source points lie in the disk of normalized radius 0.9, including the center and grid points at radii below 0.3. The inner radius 0.3 defines the initial annular illumination profile; it does not restrict the source weights during optimization. The weights satisfy (\ref{eq:source}), so the modeled source has unit total flux. The frozen optical configuration uses NA 1.35, wavelength 193\,nm, a 4\,nm pixel pitch, and a $128\times128$ raster.

The source is the only optimized variable. Four original FIT layouts and three fixed synthetic FIT layouts form the seven-layout objective. Their masks and targets are held constant during optimization. Calibration data were excluded from optimization and selection; the final3 layouts were never indexed or evaluated.

\subsection{Optimization and Qualification}
The registered solver is a Frank--Wolfe procedure with linear minimization over the guarded source-weight domain. Each seed has 68 steps at each of $\beta=200,400,800$, for 204 scheduled steps. Hard metrics are recorded at registered checkpoints. The objective weights each of the seven layouts equally throughout; there are no dynamic or hotspot-dependent layout weights.

A checkpoint can qualify only if it preserves the original FIT limits (mean PV-band area at most 234.5 pixels, mean nominal $L_2$ equal to zero, and mean worst-dose $L_2$ at most 150.25 pixels), strictly improves the new-three mean PV-band relative to the fixed seed-17 reference, does not worsen the new-three nominal or worst-dose $L_2$ means, has no blank positive target at any tested dose, and passes the registered nominal-polytope and critical-guard audits. Calibration is opened only after all five seeds yield complete, qualified FIT vectors; it is not an optimizer input or tie-breaker.

\section{Experimental Setup}
The primary evidence is the frozen October 9 schema-6 coverage run recorded in \texttt{schema6\_a7f5104\_report.json}; its full SHA-256 is listed in the accompanying README. The registered seeds were 17, 29, 43, 71, and 101. Each seed obtained a distinct feasible start from the registered initializer, passed the nominal-polytope and critical-guard audits, and completed all 204 scheduled steps with four per-layout diagnostic snapshots at the registered beta checkpoints. No fallback was used. The minimum pairwise $L_1$ distance among the five starts was 0.374521. Each run recorded three accepted updates in total (one at each beta) and 201 stationary steps. All five runs reached the same final FIT source.

The report status is \texttt{no\_fit\_qualified\_checkpoint}; every seed had zero qualified checkpoints. Total runtime was 642.929195 seconds. Calibration remained closed, and the final3 layouts were never indexed or evaluated. The run therefore produced no qualified source vector and supplies no calibration or final3 result. The earlier October 7 schema-5 fallback attempt is historical context and is not treated as an independent replicate of the October 9 run.

\section{Results}
Table~\ref{tab:reference} reports the original-FIT reference metrics and compares the seed-17 reference with the last FIT checkpoint on the same new-three layouts. The last checkpoint ties the new-three PV-band and worst-dose means and has one additional aggregate nominal mismatch.

\begin{table*}[t]
\caption{Reference and last FIT-checkpoint metrics. Values are mean pixel counts over layouts in each group.}
\label{tab:reference}
\centering
\begin{tabular}{lrrr}
\hline
FIT evaluation & PV-band & Nominal $L_2$ & Worst-dose $L_2$ \\
\hline
Original four FIT reference & 234.500 & 0.000 & 150.250 \\
New three FIT reference & 161.333 & 35.333 & 93.333 \\
New three FIT, last checkpoint & 161.333 & 35.667 & 93.333 \\
\hline
\end{tabular}
\end{table*}

No optimized checkpoint passed the registered qualification gates. All five distinct starts reached the same final FIT source. Their last FIT checkpoints tied the reference PV-band and worst-dose $L_2$ means but had nominal $L_2$ mean 35.667 rather than 35.333. Across the three new layouts, that mean difference is one additional aggregate nominal mismatch. The result failed the new FIT gate: the PV-band did not strictly improve and nominal $L_2$ worsened. The five runs are repeated starts on the same fixed FIT layouts, not independent evidence of generalization, and a stationary Frank--Wolfe gap is not a certificate of global optimality. Calibration remained closed, and the final3 layouts were never indexed or evaluated.

The per-layout diagnostics also expose an observed surrogate tradeoff at the final $\beta=800$ phase. Between the phase boundary and the last FIT checkpoint, the mean soft loss over the original four layouts decreased by $7.372981\times10^{-6}$, while the mean over the three new layouts increased by $1.343352\times10^{-6}$. These are measured changes in the implemented differentiable surrogate for this run only. They do not establish that the seven-layout objective is generally incompatible with new-layout improvement, nor do they imply a hard-metric or generalization result.

\section{Limitations and Planned Validation}
\label{sec:limits}
This is a source-only computational pilot on seven fixed FIT layouts. Focus was held at zero while dose varied by $\pm2\%$; consequently, the run does not demonstrate a focus--dose process window. It does not evaluate mask changes, a Neural-ILT model, hotspot labels, hotspot-classification metrics, or an official LithoBench or MaskOpt test set. The differentiable softcount must not be described as a measured hard PV-band or hotspot-detection result.

No final3 result is claimed; those layouts were never indexed or evaluated.

The schema-6 protocol for distinct feasible initializations has been implemented and benchmarked in the frozen run reported here; it yielded no qualified checkpoint. In response to the observed final-phase surrogate tradeoff, a schema-7 static new-three objective ablation is now implemented and pinned in a separate prospective plan, but it has not been benchmarked. It changes the optimized soft-loss average to the three new FIT layouts while retaining the same source physics, feasible-start protocol, four original-layout hard guards, all-layout hard qualification rules, and common all-seven checkpoint ranking. No result or improvement is claimed for schema 7. Future work will separately test joint source--mask optimization with a Neural-ILT model, define any hotspot target and evaluation metric, and conduct external evaluation on official LithoBench and MaskOpt data \cite{lithobench,maskopt}. Such an evaluation will require a frozen protocol and independent test data; no generalization claim is made here.

\section{Conclusion}
We documented the implemented scope and frozen outcome of a fixed-mask, 49-weight illumination pilot. Five schema-6 runs used distinct feasible starts, completed their schedules, and reached the same final FIT source without producing a qualified checkpoint. Calibration remained closed, and the final3 layouts were never indexed or evaluated. The measured final-phase soft-loss changes motivate a schema-7 static new-three objective ablation, which remains unbenchmarked. These results support a reproducible account of the current experiment, not a claim of improved printability, hotspot detection, or generalization. Developing and evaluating a joint SMO--Neural-ILT pipeline also remains future work.

\begin{thebibliography}{00}
\bibitem{hopkins}
H. H. Hopkins, "On the diffraction theory of optical images," \textit{Proceedings of the Royal Society of London. Series A, Mathematical and Physical Sciences}, vol. 217, no. 1130, pp. 408--432, 1953, doi: 10.1098/rspa.1953.0071.

\bibitem{lithobench}
S. Zheng, H. Yang, B. Zhu, B. Yu, and M. D. F. Wong, "LithoBench: Benchmarking AI computational lithography for semiconductor manufacturing," in \textit{Advances in Neural Information Processing Systems}, vol. 36, Datasets and Benchmarks Track, 2023, doi: 10.52202/075280-1315. [Online]. Available: \url{https://proceedings.neurips.cc/paper_files/paper/2023/hash/604b9fa9e1c16284e6517d923cf9ff20-Abstract-Datasets_and_Benchmarks.html}

\bibitem{neuralilt2}
B. Jiang, L. Liu, Y. Ma, B. Yu, and E. F. Y. Young, "Neural-ILT 2.0: Migrating ILT to domain-specific and multitask-enabled neural network," \textit{IEEE Transactions on Computer-Aided Design of Integrated Circuits and Systems}, vol. 41, no. 8, pp. 2671--2684, Aug. 2022, doi: 10.1109/TCAD.2021.3109556.

\bibitem{maskopt}
Y. Hu, L. Zhuang, H. Xiang, J. Xiong, and G.-J. Nam, "MaskOpt: A large-scale mask optimization dataset to advance AI in integrated circuit manufacturing," arXiv:2512.20655, 2025. [Online]. Available: \url{https://arxiv.org/abs/2512.20655}
\end{thebibliography}

\end{document}
